\section{引言}

\subsection{文档目的}

本需求规格说明书定义“轻量级外卖服务平台”的业务边界、功能行为、业务规则、外部接口和质量要求，作为需求评审、架构设计、测试驱动开发、前后端联调和交叉验收的共同基线。文档面向课程教师、验收小组、项目经理、开发人员和测试人员。

\subsection{项目背景}

项目围绕外卖业务的四方协同场景展开。顾客需要完成浏览、选购、结算和收货；商家需要维护店铺、商品并履行订单；骑手需要承接配送任务；管理员需要完成准入与异常治理。系统在完成上述基础闭环的同时，引入 AI 客服与智能点餐、营销资产和数据看板，但不让任何外部智能服务成为基础下单和配送的单点依赖。

\subsection{产品愿景与目标}

\begin{enumerate}
  \item 建立顾客、商家、骑手和管理员职责清晰的外卖履约闭环。
  \item 以订单和配送任务的独立状态流转保证业务可理解、可追踪和可测试。
  \item 以权限受控的 AI 客服连接菜单搜索、订单查询和智能点餐，展示智能服务与传统交易系统的可控集成。
  \item 对库存、订单、营销资产和骑手接单提供后端业务校验、幂等和并发保护。
  \item 使需求、接口、数据对象、自动化测试和验收结果保持可追踪关系。
\end{enumerate}

\subsection{优先级定义}

\begin{longtable}{|L{0.16\linewidth}|L{0.30\linewidth}|L{0.44\linewidth}|}
\caption{需求优先级定义}\label{tab:priority}\\ \hline
\textbf{级别} & \textbf{含义} & \textbf{验收处理} \\ \hline
P0 & 核心基线 & 必须实现并通过正常、异常和边界验收。 \\ \hline
P1 & 计划完成的亮点 & 属于本版本基线，应实现并演示完整闭环。 \\ \hline
P2 & 可选增强 & 不作为基线验收失败条件；启用时仍须符合权限、安全和降级要求。 \\ \hline
\end{longtable}

\subsection{术语与缩略语}

\begin{longtable}{|L{0.22\linewidth}|L{0.68\linewidth}|}
\caption{术语与缩略语}\label{tab:terms}\\ \hline
\textbf{术语} & \textbf{定义} \\ \hline
有效订单 & 已支付且未取消的订单；营业额与销量默认仅统计已完成的有效订单。 \\ \hline
配送任务 & 订单进入配送阶段后生成的独立履约记录，记录骑手、任务状态和关键时间。 \\ \hline
活动任务 & 处于待接单、已接单、已取餐、配送中或异常待处理状态的配送任务。 \\ \hline
资产流水 & 钱包、积分和优惠券因业务动作产生的不可直接覆盖的变更记录。 \\ \hline
幂等 & 同一业务请求重复提交时，系统不产生重复扣款、扣库存、返还或状态跳转。 \\ \hline
AI 工具调用 & AI 客服在用户授权与后端权限校验后，对菜单搜索、订单查询等受控接口的调用。 \\ \hline
降级 & 外部 AI、地图或媒体服务不可用时，返回明确提示并切换到普通搜索、文本地址或手工操作。 \\ \hline
\end{longtable}

\subsection{参考资料}

\begin{longtable}{|L{0.34\linewidth}|L{0.18\linewidth}|L{0.38\linewidth}|}
\caption{参考资料}\label{tab:references}\\ \hline
\textbf{资料} & \textbf{提供方} & \textbf{用途} \\ \hline
《软件工程综合实践》项目要求 & 课程组 & 确定前后端分离、测试先行、RESTful 接口、需求变更和交叉验收等过程要求。 \\ \hline
轻量级外卖服务平台公开仓库 & 项目组 & 保存对外源代码、需求、接口、测试和验收材料。 \\ \hline
\end{longtable}
